
{\bf 03 (2017-2018)}

 III-4 Mesure, Rappels.

Copenhague : disparition aléatoire

Everett : refuse la disparition,

\subsection{Résultat effectif d'une mesure} % III-6
\subsubsection{Indéterminisme et aléatoire} % a)

$|g_i>$

$\sum_i\alpha_i|g_i>$

\[
\]

\subsubsection{Règle de Born} % b)

$\sum_i\alpha_i|g_i>$

$|g_i>$

$\mc{P}(g_i) \propto |\alpha_i|^2$

\[
\mc{P}(g_i)=\frac{|\alpha_i|^2}{\sum_j|\alpha_j|^2}
\]

1957

$|g_i>$ $\lambda|g_i>$

\[
\mc{P}(g_i)=\frac{|\lambda|^2}{|\lambda|^2}
\]

\[
\alpha_1|g_1>+\alpha_2|g_2>+...+\alpha_n|g_n>
\]

$\forall\lambda\in\mb{C}$

\subsection{Comment distinguer deux états} % 7
\subsubsection{Classique} % a
\subsubsection{Quantique} % b

\subsection{Préparation d'un système} % 8

\subsection{Irréversibilité} % 9

\section{Structure hermitienne, rappels} % î®
\subsection{Produit scalaire sur un $\mb{R}$-espace vectoriel} % 1
\subsubsection{Définition} % a
Forme bilinéaire symétrique définie positive
\[
\begin{array}{ c r c l }
 f : & \mc{E}\times\mc{E} & \to & \mb{R} \\
  & u,v & \mapsto & f(u,v) \\
\end{array}
\]
linéaire à droite
\[
f(u,\lambda v_1+\mu v_2)=\lambda f(u,v_1)+\mu f(u,v_2)
\]
linéaire à gauche
\[
f(\lambda u_1+\mu u_2, v)=\lambda f(u_1,v)+\mu f(u_2,v)
\]
symétrique
\[
f(u,v)=f(v,u)
\]
positive
\[
f(u,v)\geqslant 0
\]
définie
\[
f(u,u)=0_\mb{R} \Leftrightarrow u=0_\mc{E}
\]

Exemple $\mb{R}^n$
\[
\left(\begin{array}{ c }
 u_1 \\ \vdots \\ u_n \\
\end{array}\right)
\qquad
\left(\begin{array}{ c }
 u_1 \\ \vdots \\ u_n \\
\end{array}\right)
\quad\to\quad f(u,v)=u.v
\]
Norme associée
\[
||u||=\sqrt{f(u,u)}
\]

\subsubsection{Notion d'orthogonalité} % b

\[
u \perp v \quad alors\quad f(u,v)=0
\]
Soit A une partie de $\mc{E}$
\[
A^\perp = \{u\in\mc{E},\ \forall a \in A,\ f(a,u)=0\}
\]

$\to$ famille orthogonale
\[
\{u_i\}\ telle\ que\ \forall i,j,\ i\neq j \Rightarrow f(u_i,u_j)=0
\]

En dimension finie, il existe une base orthogonale
\subsubsection{Expression dans une base orthogonale} % c
\[
\{e_i\}_{i\in\{1,...,n\}} \qquad f(e_i,e_j)=\delta_{ij}
\]
\[
f(u,v)=f(\sum_iu_ie_i,\sum_jv_je_j)=\sum_i\sum_ju_iv_jf(e_i,e_j)=\sum_iu_iv_i
\]
soit
\[
\boxed{f(u,v)=\sum_iu_iv_i}
\]
